\section{Ondes de surfaces circulaires}

\subsection{Description des manipulations}

La deuxième expérience se scinde en deux parties. La première permet de déduire la tension superficielle d'un liquide. La seconde permet d'étudier visuellement les différentes interférences émergeant entre deux ondes de surface circulaires.
\begin{figure}[h!]
    \centering
    % Première figure
    \begin{minipage}[h]{0.56\linewidth}
    \includegraphics[width=\linewidth]{manip_cuve.png}
    \caption{$Z_L = 2000 \ohm$}
    \label{fig:manip_cuve}
    \end{minipage}
    \hfill
    % Deuxième figure
    \begin{minipage}[h]{0.35\linewidth}
        L'installation visible sur Fig. \ref{fig:manip_cuve} se constitue d'une bac (gris) avec un fond transparent dans lequel on peut verser un liquide. La surface du liquide est en contact d'abord avec un exitateur (dispositif pneumatique avec transmissions des oscillations accoustiques par un conduit souple) pour mesurer les longueurs d'ondes de surface. Un deuxième est nécessaire lorsque l'on souhaite mesurer des interférences.
    \end{minipage}
\end{figure}

Nous commençons par utiliser un niveau à bulle pour nous assurer que la vitre de la cuve soit la plus horizontale possible et éviter toute distorsion des ondes. Une fois la cuve correctement nivelée, nous la remplissons d'une solution d'eau savonneuse de près de 1cm. Nous placons ensuite un excitateur d’ondes sur son support, ajustons la hauteur de l’excitateur afin qu'il soit partiellement immergé et transmette par l'extrémité de la buse les oscillations générées. Cela assure que les ondes soit proprement propagées sur la surface. Enfin, nous choisissons une dizaine de fréquences d’excitation (basses fréquences) comprises entre 15 et 80 Hz que nous lisons sur un oscilloscope numérique. Les impulsions entre stroboscope (en jaune sur Fig. \ref{fig:manip_cuve}) et l'exitateur sont synchronisés au travers les différentes fréquences. Cela permet de "geler" les ondes et de mesurer précisément les longueurs d’ondes sur un écran translucide (en vert) en fonction de la fréquence générée. Les longueurs d'ondes mesurées subissent un grossissement en fonction de la diffraction sur le liquide et le miroir qu'il est nécéssaire de mesurer à partir du rapport de taille entre l'objet et sa projection sur l'écran.
\\\\
En un second temps, un deuxième exitateur est placé à une distance fixe (D = 41 $\pm 1 $mm) en veillant à ce que la longueur de tube soit identique dans les deux buses. Cela permet d'éviter un déphasage entres les deux sources. Nous mesurons ensuite au travers de l'écran les interférences générées en fonction des mêmes impulsions à basses fréquences ($\approx 25.2Hz \pm 0.5 $Hz) fournies aux deux excitateurs. Nous prenons alors une photo à ~2.5m de l'écran pour la traiter l'image et déterminer les hyperboles correspondantes. Le processus est répété à une même fréquence à D $\approx 113 \pm 1 mm$ entre les excitateurs  et une seconde fois avec à même distance en baissant la fréquence à f $\approx 19.6 \pm 0.5 Hz$.
